%==================================================================================================================
%==================================================== Objectifs ===================================================
%==================================================================================================================
\section{Objectifs}

\begin{frame}[t]{Objectifs}

  \vspace{-0.75em}

  \begin{pcbanner}{Ma contribution dans GASPER}
    \centering
    \footnotesize
      Les partenaires réalisent l'élaboration, la caractérisation et les tests des dispositifs ;
  ma thèse porte sur la \textbf{modélisation de la photocathode}.
  \end{pcbanner}

  \vspace{-0.5em}

  \begin{pcbanner}{Démarche de modélisation}
    \centering
    \figschemademarchemodelo\\[0.5em]
    \vspace{0.5em}
    \footnotesize
    \textcolor{pcNavy}{
      Partie photovoltaïque
      \hspace{1em}+\hspace{1em}
      Partie électrochimique
      \hspace{1em}$\rightarrow$\hspace{1em}
      modèle complet
    }\\
    \textcolor{pcNavy}{(en parallèle : approche par bilan détaillé, transverse aux deux parties)}

  \end{pcbanner}

  \vspace{-0.65em}

\begin{pcbanner}{Finalité}
  \centering
  \footnotesize
  Utiliser le modèle pour guider la montée en échelle (extrapolation), la conception de photocathodes
  plus performantes et permettre ainsi de limiter les essais expérimentaux.
  
\end{pcbanner}

\end{frame}

%==================================================================================================================
%============================================= Fin Objectifs ===================================================
%==================================================================================================================
